\documentclass[11pt]{article}
\usepackage{mathblatt}


\begin{document}
\pruefheftstil{Gleichungssysteme}{Prüfungsheft P10}
\noindent{\Large\bfseries Gleichungssysteme}\hfill{\small\color{mbgrau}Prüfungsheft P10}\par\medskip
\pfrueckblick
\begin{pfnr}{}{1.}{}
Stelle die Gleichung $x + y = 15$ nach $y$ um.
\par\noindent y = \leerfeld \par
\rechenplatz{2}
\fusshilfe{\mbox{1:~$15 - x$}}
\end{pfnr}
\begin{pfnr}{}{2.}{}
Zeichne die Gerade $y = 2x + 1$ in das Koordinatensystem. Nutze dazu den Achsenabschnitt und ein Steigungsdreieck.
\par\smallskip\noindent \begin{ksys}[xmin=-5,xmax=5,ymin=-5,ymax=5] \end{ksys}\par
\rechenplatz{1}
\fusshilfe{\mbox{2:~Gerade durch $(0|1)$ und $(1|3)$}}
\end{pfnr}
\begin{pfnr}{}{3.}{}
Löse die Gleichung $2x + 5 = 17$.
\rechenplatz{2}
\fusshilfe{\mbox{3:~$6$}}
\end{pfnr}
\begin{pfnr}{}{4.}{}
Multipliziere die Klammer aus. $3(x + 4)$
\par\noindent \leerfeld \par
\fusshilfe{\mbox{4:~$3x + 12$}}
\end{pfnr}
\begin{pfnr}{}{5.}{}
Berechne. $3 \cdot 2{,}40\,€$
\par\noindent \leerfeld[€] \par
\fusshilfe{\mbox{5:~$7{,}20\,€$}}
\end{pfnr}
\begin{pfnr}{}{6.}{}
Drei Hefte kosten zusammen $4{,}50\,€$. Stelle eine Gleichung für den Preis $x$ eines Hefts auf und löse sie.
\par\noindent x = \leerfeld[€] \par
\rechenplatz{3}
\fusshilfe{\mbox{6:~$1{,}50\,€$}}
\end{pfnr}
\pfstufekopf[saufstellen]{aufstellen}{in 1 der letzten 5 Prüfungen}
\pfgruppe{}
\begin{pfnr}{eigene Aufgabe}{7.}{}
Gegeben sind die Gleichungen I: $x + y = 8$ und II: $3x - y = 0$. Kreuze an, von welcher Gleichung das Zahlenpaar $(2|6)$ eine Lösung ist.
\pfkreuzzeile{nur von I}\pfkreuzzeile{nur von II}\pfkreuzzeile{von beiden}\pfkreuzzeile{von keiner}
\fusshilfe{\mbox{7:~von beiden}}
\end{pfnr}
\begin{pfnr}{eigene Aufgabe}{8.}{}
Im Schwimmbad kosten $2$ Schwimmbadkarten und $5$ Saunakarten zusammen $58\,€$. $3$ Schwimmbadkarten und $1$ Saunakarte kosten $35\,€$. Überlege, was unbekannt ist, und benenne die gesuchten Größen mit $x$ und $y$. Schreibe dann heraus, welche Zahlen rechts vom Gleichheitszeichen stehen und welche vor $x$ und $y$.
\rechenplatz{3}
\fusshilfe{\mbox{8:~vor $x$ und $y$: $2$, $5$ und $3$, $1$}}
\end{pfnr}
\pfgruppe{}
\begin{pfnr}{eigene Aufgabe}{9.}{}
Die Tabelle gehört zur Gleichung $x + y = 5$. Ergänze zu jedem $x$ den passenden Wert für $y$.
\par\smallskip\noindent \wertetabelle{x}{y}{0,1,2,3}\par
\fusshilfe{\mbox{9:~$y$: $5$, $4$, $3$, $2$}}
\end{pfnr}
\pfgruppe{}
\begin{pfnr}{eigene Aufgabe}{10.}{}
Ein Buch und ein Kalender kosten zusammen $21\,€$. Zwei Bücher und fünf Kalender kosten $69\,€$. $x$ ist der Preis eines Buchs in Euro, $y$ ist der Preis eines Kalenders in Euro. Stelle die Gleichungen I und II auf.
\par\noindent I: \leerfeld[\quad]  II: \leerfeld \par
\rechenplatz{2}
\fusshilfe{\mbox{10:~$69$}}
\end{pfnr}
\begin{pfnr}{eigene Aufgabe}{11.}{}
Eine Jugendherberge hat $23$ Zimmer. Es gibt Zweibettzimmer und Vierbettzimmer. Zusammen haben sie $68$ Betten. $x$ ist die Anzahl der Zweibettzimmer, $y$ die Anzahl der Vierbettzimmer. Stelle die Gleichungen I und II auf.
\par\noindent I: \leerfeld[\quad]  II: \leerfeld \par
\rechenplatz{2}
\fusshilfe{\mbox{11:~$68$}}
\end{pfnr}
\begin{pfnr}{eigene Aufgabe}{12.}{}
Die Summe zweier Zahlen ist $31$. Ihre Differenz ist $7$. $x$ ist die größere Zahl, $y$ die kleinere Zahl. Stelle die Gleichungen I und II auf.
\par\noindent I: \leerfeld[\quad]  II: \leerfeld \par
\rechenplatz{2}
\fusshilfe{\mbox{12:~$7$}}
\end{pfnr}
\begin{pfnr}{eigene Aufgabe}{13.}{}
Die Summe zweier Zahlen ist $42$. Die größere Zahl ist das Doppelte der kleineren. Welche Zahlen sind es? Stelle ein Gleichungssystem auf, löse es und schreibe einen Antwortsatz.
\par\noindent \leerfeld[und]  \leerfeld \par
\rechenplatz{3}
\fusshilfe{\mbox{13:~$28$}}
\end{pfnr}
\pfgruppe{}
\begin{pfnr}{eigene Aufgabe}{14.}{}
Beim Schulfest wurden $120$ Getränke verkauft. Ein Wasser kostet $0{,}80\,€$, ein Saft $1{,}20\,€$. Die Einnahmen sind $118{,}40\,€$. $x$ ist die Anzahl der verkauften Wasser, $y$ die Anzahl der Säfte. Stelle die Gleichungen I und II auf.
\par\noindent I: \leerfeld[\quad]  II: \leerfeld \par
\rechenplatz{2}
\fusshilfe{\mbox{14:~$118{,}40$}}
\end{pfnr}
\pfgruppe{}
\begin{pfnr}{P10 ’16}{\llap{\textasteriskcentered\,}15.}{4 BE}
Eine Jugendherberge hat nur Zimmer mit drei Betten und Zimmer mit fünf Betten. Die Klasse kann 16 Zimmer mit zusammen 66 Betten belegen. Bestimme, wie viele Dreibettzimmer und wie viele Fünfbettzimmer es sind.
\rechenplatz{4}
\fusshilfe{\mbox{15:~7 Drei-Bett-Zimmer; 9 Fünf-Bett-Zimmer}}
\end{pfnr}
\pfgruppe{}
\begin{pfnr}{P10 ’21}{16.}{4 BE}
Zwei Familien besuchen in Paris den Eiffelturm. Familie Beyer (zwei Erwachsene, zwei Kinder) bezahlt 77,80 € Eintritt, Familie Gouvan (ein Erwachsener, drei Kinder) 64,90 €. Stelle zuerst zwei Gleichungen zu diesen Angaben auf. Bestimme dann den Eintrittspreis für einen Erwachsenen und für ein Kind.
\par\smallskip\noindent \sachtabelle{lc}{Gleichung & Term}{(I) & \leerzelle\\ (II) & \leerzelle}\par
\rechenplatz{4}
\fusshilfe{\mbox{16:~64,90; Erwachsener 25,90 €, Kind 13,00 €}}
\end{pfnr}
\pfstufekopf[slsen]{lösen}{in 1 der letzten 5 Prüfungen}
\pfgruppe{}
\begin{pfnr}{eigene Aufgabe}{17.}{}
Gegeben sind I: $2x + 3y = 9$ und II: $5x - 3y = 14$. Kreuze an, was für die Vorzahlen bei $y$ gilt.
\pfkreuzzeile{gleiche Vorzahl}\pfkreuzzeile{Gegenzahlen}\pfkreuzzeile{weder noch}
\fusshilfe{\mbox{17:~Gegenzahlen}}
\end{pfnr}
\begin{pfnr}{eigene Aufgabe}{18.}{}
Gegeben sind I: $3x + 7y = 17$ und II: $x + 4y = 9$. Du willst eine Gleichung umstellen und in die andere einsetzen. Kreuze an, welche Gleichung dafür schon fast fertig ist.
\pfkreuzzeile{Gleichung I}\pfkreuzzeile{Gleichung II}
\fusshilfe{\mbox{18:~Gleichung II}}
\end{pfnr}
\pfgruppe{}
\begin{pfnr}{eigene Aufgabe}{19.}{}
Löse das Gleichungssystem I: $y = 2x + 1$ und II: $3x + y = 16$. Setze dazu I in II ein.
\par\noindent (\leerfeld[|]  \leerfeld[)] \par
\rechenplatz{2}
\fusshilfe{\mbox{19:~Lösung $(3|7)$}}
\end{pfnr}
\begin{pfnr}{eigene Aufgabe}{20.}{}
Löse das Gleichungssystem I: $x + y = 8$ und II: $3x - y = 4$. Addiere dazu die beiden Gleichungen.
\par\noindent (\leerfeld[|]  \leerfeld[)] \par
\rechenplatz{2}
\fusshilfe{\mbox{20:~Lösung $(3|5)$}}
\end{pfnr}
\begin{pfnr}{eigene Aufgabe}{21.}{}
Löse das Gleichungssystem I: $2x + y = 9$ und II: $4x + y = 17$. Stelle dazu I nach $y$ um und setze in II ein.
\par\noindent (\leerfeld[|]  \leerfeld[)] \par
\rechenplatz{2}
\fusshilfe{\mbox{21:~Lösung $(4|1)$}}
\end{pfnr}
\begin{pfnr}{eigene Aufgabe}{22.}{}
Löse das Gleichungssystem I: $x + 2y = 7$ und II: $x + 5y = 16$. Stelle dazu I nach $x$ um und setze in II ein.
\par\noindent (\leerfeld[|]  \leerfeld[)] \par
\rechenplatz{2}
\fusshilfe{\mbox{22:~Lösung $(1|3)$}}
\end{pfnr}
\begin{pfnr}{eigene Aufgabe}{23.}{}
Löse das Gleichungssystem I: $y = x + 2$ und II: $3x + 4y = 29$. Setze dazu I in II ein.
\par\noindent (\leerfeld[|]  \leerfeld[)] \par
\rechenplatz{2}
\fusshilfe{\mbox{23:~Lösung $(3|5)$}}
\end{pfnr}
\begin{pfnr}{eigene Aufgabe}{24.}{}
Löse das Gleichungssystem I: $y = x + 4$ und II: $5x - y = 16$. Setze dazu I in II ein.
\par\noindent (\leerfeld[|]  \leerfeld[)] \par
\rechenplatz{2}
\fusshilfe{\mbox{24:~Lösung $(5|9)$}}
\end{pfnr}
\begin{pfnr}{eigene Aufgabe}{25.}{}
Löse das Gleichungssystem I: $y = x + 5$ und II: $2x + y = -1$ durch Einsetzen.
\par\noindent (\leerfeld[|]  \leerfeld[)] \par
\rechenplatz{2}
\fusshilfe{\mbox{25:~Lösung $(-2|3)$}}
\end{pfnr}
\begin{pfnr}{eigene Aufgabe}{26.}{}
Löse das Gleichungssystem I: $y = 3x - 4$ und II: $y = x + 2$. Setze dazu die rechten Seiten gleich.
\par\noindent (\leerfeld[|]  \leerfeld[)] \par
\rechenplatz{2}
\fusshilfe{\mbox{26:~Lösung $(3|5)$}}
\end{pfnr}
\begin{pfnr}{eigene Aufgabe}{27.}{}
Löse das Gleichungssystem I: $3x + 6y = 27$ und II: $3x + y = 17$. Subtrahiere dazu die Gleichungen.
\par\noindent (\leerfeld[|]  \leerfeld[)] \par
\rechenplatz{2}
\fusshilfe{\mbox{27:~Lösung $(5|2)$}}
\end{pfnr}
\begin{pfnr}{eigene Aufgabe}{28.}{}
Löse das Gleichungssystem I: $-2x + 3y = 14$ und II: $2x + y = 2$ mit dem Additionsverfahren.
\par\noindent (\leerfeld[|]  \leerfeld[)] \par
\rechenplatz{2}
\fusshilfe{\mbox{28:~Lösung $(-1|4)$}}
\end{pfnr}
\pfgruppe{}
\begin{pfnr}{eigene Aufgabe}{29.}{}
Löse das Gleichungssystem I: $2x + y = 7$ und II: $3x - 2y = 7$. Vervielfache dazu eine Gleichung. Addiere oder subtrahiere dann die Gleichungen.
\par\noindent (\leerfeld[|]  \leerfeld[)] \par
\rechenplatz{2}
\fusshilfe{\mbox{29:~Lösung $(3|1)$}}
\end{pfnr}
\begin{pfnr}{eigene Aufgabe}{30.}{}
Löse das Gleichungssystem I: $2x + 3y = 16$ und II: $5x - 2y = 21$. Vervielfache dazu beide Gleichungen. Addiere oder subtrahiere dann die Gleichungen.
\par\noindent (\leerfeld[|]  \leerfeld[)] \par
\rechenplatz{2}
\fusshilfe{\mbox{30:~Lösung $(5|2)$}}
\end{pfnr}
\pfgruppe{}
\begin{pfnr}{eigene Aufgabe}{\llap{\textasteriskcentered\,}31.}{}
Ein Helm und ein Schloss kosten zusammen $54{,}00\,€$. Ein Verein kauft $7$ Helme und $4$ Schlösser für $316{,}20\,€$. $x$ ist der Preis eines Helms in Euro, $y$ der Preis eines Schlosses in Euro. Dann gilt I: $x + y = 54$ und II: $7x + 4y = 316{,}20$. Berechne, was ein Helm und was ein Schloss kostet.
\par\noindent Helm \leerfeld[€] , Schloss \leerfeld[€] \par
\rechenplatz{3}
\fusshilfe{\mbox{31:~$20{,}60$}}
\end{pfnr}
\pfgruppe{}
\begin{pfnr}{eigene Aufgabe}{32.}{}
Löse das Gleichungssystem I: $y - 4x = 1{,}5$ und II: $5x + 2y = 16$ mit dem Einsetzungsverfahren.
\rechenplatz{4}
\fusshilfe{\mbox{32:~Lösung $(1|5{,}5)$}}
\end{pfnr}
\pfgruppe{}
\begin{pfnr}{P10 ’22}{\llap{\textasteriskcentered\,}33.}{3 BE}
Frau Bauer will in ihrem Garten Obstbäume pflanzen. Zusammen kosten ein Apfelbaum und ein Pflaumenbaum 63,00 €. Sie kauft 9 Apfelbäume und 3 Pflaumenbäume und zahlt dafür 352,20 €. Sie schreibt dazu zwei Gleichungen auf: I $x + y = 63$; II $9x + 3y = 352{,}20$. Berechne, was ein Apfelbaum und was ein Pflaumenbaum kostet.
\rechenplatz{4}
\fusshilfe{\mbox{33:~Apfelbaum 27,20 €, Pflaumenbaum 35,80 €}}
\end{pfnr}
\pfstufekopf[svariablendeuten]{Variablen deuten}{in 1 der letzten 5 Prüfungen}
\pfgruppe{}
\begin{pfnr}{eigene Aufgabe}{34.}{}
Ein Block und ein Füller kosten zusammen $8\,€$. Vier Blöcke und zwei Füller kosten $26\,€$. Dazu gehören die Gleichungen I: $x + y = 8$ und II: $4x + 2y = 26$. Gib an, wofür $x$ und $y$ stehen.
\rechenplatz{1}
\end{pfnr}
\begin{pfnr}{eigene Aufgabe}{35.}{}
Ein Zelt und ein Schlafsack kosten zusammen $190\,€$. Ein Jugendclub kauft $4$ Zelte und $6$ Schlafsäcke für $900\,€$. Dazu gehören die Gleichungen I: $x + y = 190$ und II: $4x + 6y = 900$. Was bedeuten $x$ und $y$?
\rechenplatz{1}
\end{pfnr}
\begin{pfnr}{eigene Aufgabe}{36.}{}
Ein Imker verkauft $x$ große und $y$ kleine Gläser Honig. Es gilt $8x + 5y = 230$. Erkläre, was die Zahl $8$ und die Zahl $230$ bedeuten.
\rechenplatz{2}
\fusshilfe{\mbox{36:~$230$: Einnahmen zusammen in €}}
\end{pfnr}
\pfgruppe{}
\begin{pfnr}{P10 ’22}{37.}{2 BE}
Frau Bauer will in ihrem Garten Obstbäume pflanzen. Zusammen kosten ein Apfelbaum und ein Pflaumenbaum 63,00 €. Sie kauft 9 Apfelbäume und 3 Pflaumenbäume und zahlt dafür 352,20 €. Sie schreibt dazu zwei Gleichungen auf: I $x + y = 63$; II $9x + 3y = 352{,}20$. Gib an, wofür die Variablen $x$ und $y$ in Frau Bauers Gleichungen stehen.
\rechenplatz{1}
\end{pfnr}
\pfstufekopf[sgleichunginwortefass]{Gleichung in Worte fassen und lösen}{in 1 der letzten 5 Prüfungen}
\pfgruppe{}
\begin{pfnr}{eigene Aufgabe}{38.}{}
In einem Obstkorb liegen Äpfel und Birnen. $x$ ist die Anzahl der Äpfel, $y$ die Anzahl der Birnen. Schreibe in einem Satz, was die Gleichung $x + y = 25$ sagt.
\rechenplatz{1}
\fusshilfe{\mbox{38:~Im Korb liegen zusammen $25$ Äpfel und Birnen}}
\end{pfnr}
\begin{pfnr}{eigene Aufgabe}{39.}{}
Eine Tischtennisplatte und ein Kicker kosten zusammen $900\,€$. Eine Schule kauft $2$ Platten und $1$ Kicker für $1\,300\,€$. Dazu gehören die Gleichungen I: $x + y = 900$ und II: $2x + y = 1\,300$. Was bedeuten $x$ und $y$?
\rechenplatz{1}
\end{pfnr}
\pfpruefstein{P10 ’24}
\begin{pfnr}{}{}{}
Die Brüder Daniel und Filip kaufen ihrer Mutter zum Muttertag je einen Blumenstrauß, Daniel in Laden A, Filip in Laden B.
\end{pfnr}
\begin{pfnr}{}{a)}{2 BE}
In Laden A zahlt man für eine Rose und eine Tulpe zusammen 3,80 €. Daniel bezahlt für 6 Rosen und 5 Tulpen insgesamt 21,20 €. Stelle ein lineares Gleichungssystem zu diesen Angaben auf ($x$: Preis einer Rose, $y$: Preis einer Tulpe).
\par\smallskip\noindent \sachtabelle{lc}{Gleichung & Term}{I. & \leerzelle\\ II. & \leerzelle}\par
\rechenplatz{3}
\end{pfnr}
\begin{pfnr}{}{*b)}{4 BE}
In Laden B kostet eine Rose 2,30 € und eine Tulpe 1,70 €. Filip schreibt zu seinem Strauß das Gleichungssystem I. $2{,}30r + 1{,}70t = 26{,}90$; II. $r + t = 13$. Schreib in einem Satz auf, was Gleichung II über den Strauß aussagt. Löse das Gleichungssystem.
\rechenplatz{3}
\end{pfnr}
\end{document}